\documentclass[12pt, a4paper]{ctexart}
\usepackage{geometry}
\geometry{left=2.5cm, right=2.5cm, top=3cm, bottom=3cm}
\usepackage{amsmath}
\usepackage{circuitikz}
\usepackage{float}
\usepackage{listings}
\usetikzlibrary{arrows, shapes.gates.logic.US}
\ctikzset{
    logic ports=ieee,
    logic ports/scale=0.8,
    multipoles/flipflop/pin spacing=0.4
}
\input{risc-v.tex}
\lstset{
    basicstyle=\ttfamily\small,
    frame=single,
    framesep=5pt,
    xleftmargin=10pt,
    xrightmargin=10pt,
    breaklines=true
}

\begin{document}

\section*{第8章习题}

\subsection*{3.}

\subsubsection*{(1)}

8 位（1 字节）

\subsubsection*{(2)}

\begin{figure}[H]
    \centering
    \begin{circuitikz}
        \node[
            one bit adder,
            muxdemux def={NB=0}
        ] (adder1) at (0,0) {+};
        \node[
            one bit adder,
            muxdemux def={NB=0},
            anchor=lpin 1
        ] (adder2) at (adder1.rpin 1) {+}; 
        \node[
            muxdemux,
            muxdemux def={Lh=3, Rh=1, w=2, NL=2, NR=1, NT=0, NB=1},
            muxdemux label={L1=1, L2=0},
            anchor=lpin 2
        ] (mux) at (adder2.rpin 1) {};
        \draw (adder1.rpin 1) node[circ]{} -- ++(0,0.5) -| (mux.lpin 1);
        \node[left] at (adder1.lpin 1) {PC};
        \node[left] at (adder1.lpin 2) {2};
        \node[right] at (mux.rpin 1) {nextPC};
        \node[left] at (adder2.lpin 2) {\{8\{Imm8[7]\}, Imm8\}};
        \node[below right] at (mux.bpin 1) {
            $\overline{ZF+(SF\oplus OF)}\cdot Bgt$
        };
    \end{circuitikz}
\end{figure}

\subsection*{4.}

\begin{itemize}
    \item RegWr=0 时所有需要写入寄存器的指令无法执行
    \item ALUASrc=0 时 J 型指令无法执行
    \item Branch=0 时 B 型指令无法执行
    \item Jump=0 时 J 型指令无法执行
    \item MemWr=0 时 store 指令无法执行
    \item MemtoReg=0 时 load 指令无法执行
\end{itemize}

\subsection*{5.}

\begin{itemize}
    \item RegWr=1 时所有不需要写入寄存器的指令无法执行
    \item ALUASrc=1 时所有除 J 型指令以外需要 ALU 的指令无法执行
    \item Branch=1 时除 B 型以外的指令无法执行
    \item Jump=1 时除 J 型以外的指令无法执行
    \item MemWr=1 时除 store 以外的指令无法执行
    \item MemtoReg=1 时除 load 以外的指令无法执行
\end{itemize}

\subsection*{6.}

\subsubsection*{(1)}

\begin{lstlisting}[language={[RISC-V]Assembler}]
xor     rs, rs, rt
xor     rt, rs, rt
xor     rs, rs, rt
\end{lstlisting}

\subsubsection*{(2)}

5\%

\subsubsection*{(3)}

不能在单周期中实现，因为只有一个写入端口。多周期可以在不用时钟周期中写入，可以实现。

\subsection*{10.}

\subsubsection*{(1)}

\begin{table}[H]
    \centering
    \begin{tabular}{|c|c|c|}
        \hline
        异常或中断 & 指令类型 & 检测到的时钟周期 \\ \hline
        除数为 0 & 除法指令 & EX 阶段 \\
        溢出 & 带溢出检测的算术指令 & EX 阶段 \\
        无效指令操作码 & 非法指令 & ID 阶段 \\
        无效指令地址 & 任何指令 & IF 阶段 \\
        无效数据地址 & 访存指令 & M 阶段 \\
        缺页 & 访存指令 & M 阶段 \\
        访问越权 & 访存指令 & M 阶段 \\
        外部中断 & 任何指令 & IF 阶段 \\ \hline
    \end{tabular}
\end{table}

\subsubsection*{(2)}

\begin{enumerate}
    \item 保存断点和程序状态：正确计算断点值，程序状态储存在 PSWR 中。
    \item 识别异常/中断源并转相应处理程序执行：使用软件识别内部异常，使用软件或硬件识别外部中断。
\end{enumerate}

\subsection*{11.}

\subsubsection*{(1)}

不能，因为最慢的还是存储单元的 200ps。

\subsubsection*{(2)}

无影响，因为最慢的还是存储单元的 200ps。

\subsubsection*{(3)}

有影响，现在最慢的是 ALU 的 210ps，会让流水线的时钟周期变为 210ps。

\subsection*{13.}

\subsubsection*{(1)}

\begin{table}[H]
    \centering
    \begin{tabular}{|ccc|ccc|}
        \hline
        80ps & 30ps & 60ps & 50ps & 70ps & 10ps \\ \hline
        \multicolumn{3}{|c|}{170ps} & \multicolumn{3}{|c|}{130ps} \\ \hline
    \end{tabular}
\end{table}

如图，时钟周期为 190ps，指令吞吐率为 5.26 GIPS，指令执行时间为 380ps。

\subsubsection*{(2)}

\begin{table}[H]
    \centering
    \begin{tabular}{|cc|cc|cc|}
        \hline
        80ps & 30ps & 60ps & 50ps & 70ps & 10ps \\ \hline
        \multicolumn{2}{|c|}{110ps} & \multicolumn{2}{|c|}{110ps} &
        \multicolumn{2}{|c|}{80ps} \\ \hline
    \end{tabular}
\end{table}

如图，时钟周期为 130ps，指令吞吐率为 7.69 GIPS，指令执行时间为 390ps。

\subsubsection*{(3)}

\begin{table}[H]
    \centering
    \begin{tabular}{|c|cc|c|cc|}
        \hline
        80ps & 30ps & 60ps & 50ps & 70ps & 10ps \\ \hline
        \multicolumn{1}{|c|}{80ps} & \multicolumn{2}{|c|}{90ps} &
        \multicolumn{1}{|c|}{50ps} & \multicolumn{2}{|c|}{80ps} \\ \hline
    \end{tabular}
\end{table}

如图，时钟周期为 110ps，指令吞吐率为 9.09 GIPS，指令执行时间为 440ps。

\subsubsection*{(4)}

\begin{table}[H]
    \centering
    \begin{tabular}{|c|c|c|c|cc|}
        \hline
        80ps & 30ps & 60ps & 50ps & 70ps & 10ps \\ \hline
        \multicolumn{1}{|c|}{80ps} & \multicolumn{1}{|c|}{30ps} &
        \multicolumn{1}{|c|}{60ps} & \multicolumn{1}{|c|}{50ps} &
        \multicolumn{2}{|c|}{80ps} \\ \hline
    \end{tabular}
\end{table}

如图，时钟周期为 100ps，指令吞吐率为 10.00 GIPS，指令执行时间为 500ps。

\subsection*{14.}

\subsubsection*{(1)}

(1, 2), (2, 3), (2, 4), (3, 4)

\subsubsection*{(2)}

假设寄存器堆支持半周期读写，修改后如下

\begin{lstlisting}[language={[RISC-V]Assembler}]
add     s3, s1, s0
nop
nop
add     t2, s3, s3
nop
nop
lw      t1, 0(t2)
nop
nop
add     t3, t1, t2
\end{lstlisting}

\subsubsection*{(3)}

不能完全解决，修改后如下

\begin{lstlisting}[language={[RISC-V]Assembler}]
add     s3, s1, s0
add     t2, s3, s3
lw      t1, 0(t2)
nop
add     t3, t1, t2
\end{lstlisting}

\subsection*{15.}

\subsubsection*{(1)}

\begin{itemize}
    \item 从指令 1 的 EX/M 转发到指令 2 的 EX 阶段
    \item 从指令 2 的 EX/M 转发到指令 3 的 EX 阶段
    \item 从指令 3 的 M/WB 转发到指令 4 的 EX 阶段
\end{itemize}

\subsubsection*{(2)}

存在于指令 3 和指令 4 之间

\subsubsection*{(3)}

\begin{itemize}
    \item 指令 1 在 WB 阶段
    \item 指令 1 在 M 阶段
    \item 指令 1 在 EX 阶段
    \item 指令 1 在 ID 阶段
    \item 指令 1 在 IF 阶段
\end{itemize}

\begin{itemize}
    \item t1, t2 正被读出
    \item s3 将被写入
\end{itemize}

\subsection*{17.}

\begin{lstlisting}[language={[RISC-V]Assembler}]
lw      s2, 100(s6)
add     s6, s4, s7
add     s2, s2, s3
lw      s3, 200(s7)
lw      s2, 300(s8)
sub     s3, s4, s6
beq     s2, s8, Loop
\end{lstlisting}

\subsection*{18.}

\subsubsection*{(1)}

2 个时钟周期

\subsubsection*{(2)}

\texttt{bne t0, s5, Exit} 会发生阻塞，需要阻塞 1 个时钟周期

\end{document}
